\documentclass[12pt]{article}
\usepackage{amsmath, amssymb}
\usepackage[margin=1in]{geometry}
\setlength{\parskip}{6pt}
\title{QUBO Formulation: Minimum Missing Edges Hamiltonian Path}
\date{}
\begin{document}
\maketitle

\subsection*{Minimum Missing Edges Hamiltonian Path}

\paragraph{Problem Statement.}
Given an undirected graph $G=(V,E)$ with $|V|=N$ vertices and adjacency matrix $A \in \{0,1\}^{N \times N}$, find a permutation of the vertices that forms a path visiting each vertex exactly once, while minimizing the number of consecutive pairs in the path that are not connected by an edge in $G$.

\paragraph{Decision Variables.}
Introduce binary variables $x_{i,v} \in \{0,1\}$ for each position index $i \in \{0, 1, \dots, N-1\}$ and each vertex $v \in V$. The variable $x_{i,v}$ equals $1$ if and only if vertex $v$ is assigned to position $i$ in the path; otherwise, it equals $0$.

\paragraph{QUBO Derivation.}
\textbf{Step 1 --- Objective Function.} The number of missing edges between consecutive positions $i$ and $i+1$ in the path is given by $\sum_{u \in V} \sum_{v \in V} (1 - A_{uv}) x_{i,u} x_{i+1,v}$. Summing over all consecutive position pairs yields the objective to be minimized:
\begin{align}
H_{\text{obj}}(x) = \sum_{i=0}^{N-2} \sum_{u \in V} \sum_{v \in V} (1 - A_{uv}) x_{i,u} x_{i+1,v}.
\end{align}

\textbf{Step 2 --- Constraints.} The assignment must satisfy two hard constraints. First, each position must contain exactly one vertex:
\begin{align}
\sum_{v \in V} x_{i,v} = 1, \quad \forall i \in \{0, \dots, N-1\}.
\end{align}
Second, each vertex must be used exactly once in the path:
\begin{align}
\sum_{i=0}^{N-1} x_{i,v} = 1, \quad \forall v \in V.
\end{align}

\textbf{Step 3 --- Penalty Term Conversion.} We convert each constraint into a quadratic penalty term with weight $\lambda > 0$. For the position assignment constraint, we write:
\begin{align}
P_1(x) = \lambda \sum_{i=0}^{N-1} \left( \sum_{v \in V} x_{i,v} - 1 \right)^2.
\end{align}
Expanding the square and applying the idempotent property $x_{i,v}^2 = x_{i,v}$ for binary variables:
\begin{align}
\left( \sum_{v \in V} x_{i,v} - 1 \right)^2 &= \sum_{v \in V} x_{i,v}^2 + 2 \sum_{v < v' \in V} x_{i,v} x_{i,v'} - 2 \sum_{v \in V} x_{i,v} + 1 \\
&= \sum_{v \in V} x_{i,v} + 2 \sum_{v < v' \in V} x_{i,v} x_{i,v'} - 2 \sum_{v \in V} x_{i,v} + 1 \\
&= 1 - \sum_{v \in V} x_{i,v} + 2 \sum_{v < v' \in V} x_{i,v} x_{i,v'}.
\end{align}
Substituting this back into $P_1(x)$ gives:
\begin{align}
P_1(x) = \lambda \sum_{i=0}^{N-1} \left( 1 - \sum_{v \in V} x_{i,v} + 2 \sum_{v < v' \in V} x_{i,v} x_{i,v'} \right).
\end{align}
Similarly, for the vertex usage constraint:
\begin{align}
P_2(x) = \lambda \sum_{v \in V} \left( 1 - \sum_{i=0}^{N-1} x_{i,v} + 2 \sum_{i < i'} x_{i,v} x_{i',v} \right).
\end{align}

\textbf{Step 4 --- Combined Hamiltonian.} The total QUBO Hamiltonian is the sum of the objective and penalty terms:
\begin{align}
H(x) = H_{\text{obj}}(x) + P_1(x) + P_2(x).
\end{align}

\textbf{Step 5 --- Q Matrix and Offset.} Mapping the variables to a single index $k = iN + v$, the Hamiltonian takes the standard QUBO form $H(x) = \sum_{k} Q_{k,k} x_k + \sum_{k < l} Q_{k,l} x_k x_l + c$. Reading off the coefficients from the expanded expressions yields the following closed-form entries for the symmetric upper-triangular Q matrix and the constant offset $c$:
\begin{align}
Q_{(i,u),(j,v)} &= 
\begin{cases} 
-2\lambda & \text{if } i=j \text{ and } u=v, \\
2\lambda & \text{if } i=j \text{ and } u \neq v, \\
2\lambda & \text{if } u=v \text{ and } i \neq j, \\
1 - A_{uv} & \text{if } j = i+1 \text{ and } u \neq v, \\
1 - A_{uv} + 2\lambda & \text{if } j = i+1 \text{ and } u = v, \\
0 & \text{otherwise},
\end{cases} \\
c &= 2N\lambda.
\end{align}

\paragraph{Penalty Weight Justification.}
The penalty weight $\lambda$ must be chosen large enough to ensure that any violation of the constraints increases the total energy more than the maximum possible reduction in the objective function. The objective coefficients are bounded by $\max_{u,v} |1 - A_{uv}| = 1$. A single constraint violation can alter at most $N$ variable assignments, affecting at most $N-1$ consecutive position pairs in the path. Thus, the maximum objective gain from infeasibility is bounded by $N-1$. Choosing $\lambda > N-1$ guarantees that the penalty energy strictly dominates any feasible objective improvement, forcing the optimizer toward valid permutations.

\paragraph{Correctness Argument.}
Minimizing $H(x)$ over $\{0,1\}^{N^2}$ recovers the optimal Hamiltonian path because (i) any assignment satisfying both constraints corresponds to a valid permutation and incurs zero penalty energy, reducing $H(x)$ exactly to $H_{\text{obj}}(x)$; (ii) any infeasible assignment incurs a penalty energy of at least $\lambda$, which exceeds the maximum possible reduction in the objective term achievable by violating constraints. Consequently, the global minimum of $H(x)$ must occur at a feasible point, and among feasible points, the minimizer corresponds to the permutation with the minimum number of missing consecutive edges.

\paragraph{Illustrative Example.}
Consider a concrete instance with $N=3$ vertices $V=\{0,1,2\}$ and edges $E=\{(0,1), (1,2)\}$, forming a path graph. The adjacency matrix entries are $A_{01}=A_{10}=A_{12}=A_{21}=1$, and all other entries are $0$. Setting $\lambda=10$, the offset becomes $c = 2(3)(10) = 60$. The concrete objective is $H_{\text{obj}}(x) = \sum_{i=0}^{1} \sum_{u,v \in \{0,1,2\}} (1 - A_{uv}) x_{i,u} x_{i+1,v}$. The concrete penalty terms are $P_1(x) = 10 \sum_{i=0}^{2} (\sum_{v=0}^{2} x_{i,v} - 1)^2$ and $P_2(x) = 10 \sum_{v=0}^{2} (\sum_{i=0}^{2} x_{i,v} - 1)^2$. The resulting Hamiltonian is $H(x) = H_{\text{obj}}(x) + P_1(x) + P_2(x)$ with offset $c=60$. Instantiating the Q matrix entries: all linear terms are $Q_{(i,v),(i,v)} = -20$; the pair at positions $0$ and $1$ with vertices $0$ and $1$ (indices $k=0$ and $k=4$) has coefficient $Q_{(0,0),(1,1)} = 1 - A_{01} = 0$; the pair at positions $0$ and $1$ with vertices $0$ and $0$ (indices $k=0$ and $k=3$) has coefficient $Q_{(0,0),(1,0)} = 1 - A_{00} + 20 = 20$; the pair at positions $0$ and $0$ with vertices $0$ and $1$ (indices $k=0$ and $k=1$) has coefficient $Q_{(0,0),(0,1)} = 20$. Evaluating $H(x)$ on the valid permutation $0 \to 1 \to 2$ yields an objective value of $0$, zero penalty energy, and a total energy of $60$. Any permutation violating the constraints incurs an energy of at least $60 + 10 = 70$, confirming that the global minimum correctly identifies the optimal path.

\end{document}